\documentclass[11pt]{letter}
\usepackage[a4paper,margin=2.5cm]{geometry}
\signature{Salvatore Capuozzo\\University of Naples Federico II\\salvatore.capuozzo@unina.it}
\address{[Department]\\University of Naples Federico II\\Via Claudio 21, 80125 Naples, Italy}
\begin{document}
\begin{letter}{Editor-in-Chief and Guest Editors\\Neural Computing and Applications\\Topical Collection ``Emerging Trends in Smart Waste Monitoring and Environmental Surveillance''}
\opening{Dear Editors,}

we submit the manuscript ``TrashScribe: Segmenting and Describing Waste across Environments on an Edge Device'' for consideration in the Topical Collection \emph{Emerging Trends in Smart Waste Monitoring and Environmental Surveillance}.

The manuscript addresses three gaps named in the call: weak detection across varied environments, limited and incompatible datasets, real-time operation on resource-constrained devices. It defines a cross-environment benchmark that unifies five public waste datasets (street, recycling conveyor, aerial, indoor, beach) and tests each environment left out of training. It proposes \emph{TrashScribe}, a 34\,M-parameter network that segments waste and describes every object by its physical properties (material, rigidity, optical behaviour, shape) in one pass on a Jetson Orin Nano. It then measures when this small model can replace the 848\,M-parameter foundation model SAM~3. Trained with an environment, TrashScribe beats SAM~3 on all five; on a new environment ten labelled images, SAM~3's own labels on unlabelled images or, indoors, a few hundred generic photographs without waste bring it to or above SAM~3 on four of the five, while the complete system runs at about five frames per second at 15\,W.

\textbf{Relation to previous work.} The manuscript is original, has not been published and is not under consideration elsewhere. It builds on two works by the authors that are cited anonymously for double-blind review: [LENTE reference and status], a parameter-efficient SAM2 adaptation for landslide segmentation whose adapter recipe we reuse for a different domain and task; and [``Describing, Not Naming'' reference and status], a descriptor-based retrieval system for parasite-egg microscopy whose descriptor-head principle we transfer from object crops to dense segmentation. Problem, data, architecture components (dense descriptor field, intrinsic/contextual pooling, cross-dataset transplant, environment adversary, edge pipeline), experiments and text are new; no text, table or figure is reused.

\textbf{Possible conflicts.} [Two Guest Editors of this Topical Collection are affiliated with the University of Naples Federico II, as are the authors. We kindly ask that the manuscript be handled by an editor without this affiliation.]

\textbf{Suggested reviewers} (institutional e-mail, no conflicts): [Name, affiliation, e-mail] $\times$ 3--5.

All authors approved the submission. The code and configuration files are available to reviewers at [anonymised URL].

\closing{Sincerely,}
\end{letter}
\end{document}
